\documentclass[11pt,a4paper]{article}
\usepackage[T1]{fontenc}
\usepackage{newtxtext}
\usepackage{sourcesans}
\usepackage{newtxmath}
\usepackage[margin=24mm,top=21mm,bottom=20mm]{geometry}
\usepackage{booktabs,tabularx,xcolor,microtype}
\usepackage[colorlinks=true]{hyperref}

\definecolor{ink}{HTML}{242A26}
\definecolor{accent}{HTML}{2D5748}
\definecolor{muted}{HTML}{505D55}
\hypersetup{
  urlcolor=accent,
  pdftitle={Forecasting the effect of punishment on group earnings},
  pdfauthor={Ashmit Mukherjee},
  pdfsubject={A retrospective evaluation across 20 public-goods game designs}
}
\pagestyle{empty}
\setlength{\parindent}{0pt}
\setlength{\parskip}{6pt}
\newcommand{\minorheading}[1]{%
  \par\vspace{5pt}{\sffamily\bfseries\large\color{accent}#1}\par\nobreak}
\newcommand{\advisor}[2]{\href{#2}{#1}}

\begin{document}
\color{ink}
{\sffamily\small\color{muted} NYU Abu Dhabi \hfill Computer Science Capstone \enspace / \enspace November 2026}\par
{\color{accent}\hrule height .6pt}
\vspace{15pt}
{\fontsize{26}{29}\selectfont Forecasting the effect of punishment\\on group earnings\par}
\vspace{2pt}
{\sffamily\bfseries Ashmit Mukherjee}\par
{\sffamily\small Advised by Professor \advisor{Hanan Salam}{https://scholar.google.com/citations?user=9lTiKjwAAAAJ} and Professor \advisor{Benjamin Rosche}{https://scholar.google.com/citations?user=GAxIreYAAAAJ}}\par

\minorheading{The forecasting task}
Punishment can encourage people to contribute more to a common fund, but imposing a penalty costs the group money. Can a language model predict the net effect on earnings? I used the 20 public-goods game designs that Alsobay and colleagues gave to human forecasters in their validation study.

Each model saw 15 settings for one game, including group size, the return on contributions, and the cost of punishment. I gave it the average earnings of human groups that played without punishment and asked it to predict earnings with punishment. Earnings are expressed as a percentage of what the group would earn if everyone always contributed everything.

I fixed the prompt and scoring rules before collecting model forecasts. The human outcomes were already public. I averaged each model’s valid forecasts within each design, then calculated root mean squared error (RMSE) across the 20 designs, weighting each design equally.

\vspace{6pt}
\noindent\begin{minipage}[t]{.48\textwidth}\vspace{0pt}
{\sffamily\bfseries\large\color{accent}Forecast error}\par
{\small Lower RMSE indicates greater accuracy.}\par
\vspace{3pt}
{\small\renewcommand{\arraystretch}{1.18}%
\begin{tabularx}{\linewidth}{@{}Xr@{}}
\toprule
\textbf{Forecast} & \textbf{RMSE (points)}\\
\midrule
Elastic net & 4.52\\
Ordinary linear regression & 5.22\\
Claude Haiku 4.5 & 6.21\\
Lay crowd & 6.42\\
Expert crowd & 6.55\\
Claude Sonnet 4.5 & 6.62\\
Claude Sonnet 5.5 & 6.65\\
Predict no change & 7.14\\
GPT-5 & 9.28\\
\bottomrule
\end{tabularx}}\par
{\footnotesize\color{muted} Both regressions are reported in the source study. I recalculated the crowd scores from its archived forecasts.}\par
\end{minipage}\hfill
\begin{minipage}[t]{.47\textwidth}\vspace{0pt}
\setlength{\parskip}{5pt}
{\sffamily\bfseries\large\color{accent}What the forecasts missed}\par
\raggedright
Across these 20 designs, group earnings fell by 2.2 percentage points on average with punishment. Three models and the expert crowd expected a gain. Sonnet 5.5 alone among the models predicted a loss. Its training cutoff came after publication of the human outcomes.

The three Claude models and both human crowds improved on predicting no change by less than one point of RMSE. The elastic net improved on it by 2.62 points. GPT-5 had higher error than the no-change forecast and provided fewer forecasts per design after reaching its spending limit.

Haiku had the lowest error among the language models. The 95\% interval for its RMSE minus the experts' was $-0.89$ to $+0.17$ points, from 10,000 paired resamples of the designs.
\end{minipage}

\minorheading{From forecasting to playing}
I built a simulator for the seven designs played without chat. A two-design pilot tested whether groups of model players could finish a game. Whether playing improves on a direct forecast remains open.

\vfill
{\color{accent}\hrule height .5pt}
\vspace{4pt}
{\sffamily\small\color{muted}
\href{https://arxiv.org/abs/2508.17151}{\textit{Integrative Experiments Identify How Punishment Impacts Welfare in Public Goods Games}} (2025).\\
\href{https://scholar.google.com/citations?user=DxlsyhcAAAAJ}{Mohammed Alsobay},
\href{https://scholar.google.com/citations?user=C0ANojIAAAAJ}{David G. Rand},
\href{https://scholar.google.com/citations?user=LhOAiXMAAAAJ}{Duncan J. Watts}, and
\href{https://scholar.google.com/citations?user=nrEeeY4AAAAJ}{Abdullah Almaatouq}.
\enspace\href{https://doi.org/10.17605/OSF.IO/2D56W}{Data archive.}}\par
\end{document}
